\subsection*{39. Cash Reserves Need a Purpose, Not a Fashionable Multiple}

Founders sometimes copy an arbitrary rule such as ``keep 6 months of expenses.'' A fixed multiple can be a useful starting heuristic, but it is not scientifically optimal for every business. A prepaid consumer product with short delivery cycles, low fixed costs, and reliable demand has different risks from a laboratory with long procurement lead times and concentrated funding.

A more defensible reserve policy considers:
\begin{itemize}
\item Unavoidable cash commitments,
\item Volatility of collections,
\item Time needed to reduce costs responsibly,
\item Supplier and payroll obligations,
\item Customer refunds and delivery exposure,
\item Debt maturities, and
\item Time required to secure alternative funding.
\end{itemize}
Include a cushion for estimation error and major low-frequency disruptions.

Partition cash conceptually into 4 distinct buckets:
\begin{enumerate}
\item \textbf{Operating float} for normal timing variation;
\item \textbf{Obligation cover} for prepaid and legally owed services;
\item \textbf{Contingency reserve} for unexpected operational shocks; and
\item \textbf{Discretionary investment capital} for validated, reversible growth.
\end{enumerate}
Avoid double-counting the same liability in multiple categories, but do not spend the same rupee twice on competing plans.

\subsection*{40. Reserves and the Cost of Waiting}

Holding reserves has an economic cost: money could be invested in product, distribution, or equipment. But it also has option value: it allows a founder to defer hiring until a real bottleneck emerges, negotiate without desperation, and survive a delayed customer payment without breaking trust.

A thoughtful reserve policy balances those forces. The company should be willing to spend above the minimum reserve on a sufficiently attractive, \emph{reversible or staged} experiment. It should be much more conservative about an irreversible obligation such as a long lease, debt secured by personal property, or permanent headcount unsupported by repeatable revenue.

The relevant question is not ``Are we courageous enough to spend the cash?'' It is ``What additional evidence does spending this cash buy, how large is the irreversible downside, and how many future experiments will be lost if the answer is negative?''

\subsection*{41. Debt Changes the Distribution of Outcomes}

Borrowing can finance productive assets and smooth a legitimate timing mismatch. It can also convert a business-model hypothesis into an inflexible payment obligation. The distinction is particularly important when debt service begins before customer receipts become reliable.

The appropriate assessment includes interest, fees, amortization schedule, collateral, guarantees, covenants, refinancing requirements, variable-rate exposure, and the consequences of missed payments. An interest-only grace period is not evidence that principal can safely be repaid later. A short-term facility used to finance a long-lived speculative asset creates a dangerous maturity mismatch.

Debt is most defensible when repayment comes from a credible and appropriately stressed operating cash stream or from an asset with sound collateral and residual value; not from perpetual assumptions that future sales will appear on demand. Avoid relying on customer deposits to service unrelated personal debt: that may force the venture to compromise its promises at the worst moment.

\subsection*{42. Debt-Service Coverage and Covenant Headroom}

For lending contexts, a \textbf{debt-service coverage ratio (DSCR)} compares defined cash earnings with scheduled interest and principal payments:
\begin{equation}
\text{DSCR} = \frac{\text{Cash for Debt Service}}{\text{Scheduled Debt Service}}.
\end{equation}
Definitions vary materially by lender and contract. A ratio above 1 suggests more earnings or cash flow under that definition than scheduled service, but it does not establish that all payments arrive on time or that the company has adequate cash for capital expenditure and working capital.

Suppose a fictional lender uses cash available for debt service of ₹15 lakh annually and scheduled annual service is ₹12 lakh. The illustrative ratio is:
\begin{equation}
\text{DSCR} = \frac{₹15,00,000}{₹12,00,000} = \mathbf{1.25}.
\end{equation}
If collections slow or operating costs rise, the cushion may vanish. If the loan covenant requires a particular definition, the covenant definition; not the founder's preferred spreadsheet metric; is what matters.

Maintain a separate covenant register and forecast. Test whether adverse scenarios cause breaches before they occur. A company can have enough cash for this month's bills yet violate a lender's covenant that restricts future borrowing or accelerates obligations.

\subsection*{43. Financing Readiness Begins Before Funding Is Urgent}

The time to prepare a financing plan is before the company faces a deadline. A lender or investor will want evidence of demand, financial records, coherent assumptions, contracts, collections, debt exposures, and a defensible use of funds. An emergency raise may force a founder to accept poor terms, surrender valuable control, or divert attention from the product.

Map financing options in terms of \textbf{availability}, \textbf{cost}, \textbf{security}, \textbf{recourse}, \textbf{dilution}, \textbf{restrictions}, and \textbf{time to cash}. Founder savings, retained profits, customer prepayments, supplier credit, bank lending, revenue-based arrangements, strategic contracts, grants, and equity have very different burdens. A source that is theoretically available in the market may not be available to \emph{this venture today}.

Do not treat a loan approval, investor term sheet, or grant application as collected cash. Keep the base-case solvency plan viable without speculative financing. Then decide which financing opportunities accelerate a proven economics rather than substitute for it.

\subsection*{44. Return on Invested Capital Is Meaningless If the Firm Cannot Bridge the Timing Gap}

A founder can identify an apparently high-return project and still be unable to afford it. Imagine a service contract that yields ₹4 lakh of economic contribution over 9 months but requires ₹7 lakh upfront from the company. If available cash is ₹3 lakh and the buyer provides no deposit, the opportunity is not executable without outside financing or restructured terms.

It may be rational to negotiate staged milestones, reduce the initial scope, subcontract differently, or decline. The correct decision is not determined by the project margin alone. It depends on the \emph{funding gap}, reliability of collections, customer concentration, availability of capital, and downside obligations if the project is delayed.

Good capital allocation asks both whether a project creates value and whether financing it leaves the rest of the company intact. A valuable opportunity that exposes every other customer to nonperformance may be strategically unattractive.
